\section*{第一单元 负数}

\begin{definition}[负数]
    比 0 小的数叫做负数（0 除外）。负数都比自然数小。负数用负号“-”标记，如-2, -5.33, -45, -0.6 等。
\end{definition}

\begin{definition}[正数]
    大于 0 的数叫正数（不包括 0），数轴上 0 右边的数叫做正数。若一个数大于零（$>0$），则称它是一个正数。正数的前面可以加上正号“+”来表示。正数有无数个，其中有正整数，正分数和正小数。
\end{definition}

\begin{definition}[0]
    0 既不是正数，也不是负数，它是正、负数的分界数。正数都大于 0，负数都小于 0，正数大于一切负数。
\end{definition}

\begin{theorem}[在直线上表示数]
    \begin{enumerate}
        \item 正数、0 和负数可以用直线上的点表示出来。直线上的每一个点都与一个数相对应，任何一个数都可以用直线上的点来表示。
        \item 用有正数和负数的直线可以表示距离和相反的方向。
    \end{enumerate}
\end{theorem}

\subsection*{题型}
\begin{example}
    将以下数字按要求分类：
    $1.25$、$ \frac{3}{5}$、$-7$、$3$、$3.011\dots$、$-5\frac{1}{2}$、$0$、$\frac{12}{7}$、$-0.03$
\end{example}

\v{2em}

\begin{solution}
    \begin{itemize}
        \item \textbf{正数}: $1.25, \frac{3}{5}, 3, 3.011\dots, \frac{12}{7}$
        \item \textbf{负数}: $-7, -5\frac{1}{2}, -0.03$
        \item \textbf{自然数}: $3, 0$
        \item \textbf{非正数}: $-7, -5\frac{1}{2}, 0, -0.03$
    \end{itemize}
\end{solution}

\begin{example}
    写出下列数相对的负数形式：
    $0.33\dots$、$\frac{7}{19}$、$3\frac{1}{3}$、$2$、$+\frac{7}{5}$、$ +3$
\end{example}

\v{2em}

\begin{solution}
    $-0.33\dots$、$-\frac{7}{19}$、$-3\frac{1}{3}$、$-2$、$-\frac{7}{5}$、$-3$
\end{solution}

\begin{example}
    如果 $+20\%$ 表示增加 $20\%$，那么 $-20\%$ 表示什么？
\end{example}

\v{2em}

\begin{solution}
    减少 $20\%$。
\end{solution}

\begin{example}
    某日傍晚，黄山的气温由上午的零上 2 摄氏度下降了 7 摄氏度，这天傍晚黄山的气温是\underline{\hspace{1cm}}摄氏度。
\end{example}

\v{2em}

\begin{solution}
    $2 - 7 = -5$。这天傍晚黄山的气温是 \textbf{-5} 摄氏度。
\end{solution}

\begin{example}
    在数轴上表示下列个数：$1.75, -3, \frac{1}{4}, 5, 0, -3.2, -4\frac{3}{4}$
\end{example}

\v{2em}

\begin{solution}
    (此题需要在数轴上画点，此处省略图示)
\end{solution}

\begin{example}
    写出下列各点表示的数

    (示意图：A B C D E F G 分别在 -7, -5, -2, -1, 1, 4, 8 的位置)
\end{example}

\v{2em}

\begin{solution}
    A: -7, B: -5, C: -2, D: -1, E: 1, F: 4, G: 8
\end{solution}


\section*{第二单元 百分数（二）}

\begin{definition}[折扣]
    几折就是十分之几，也就是百分之几十。例如：八五折表示现价是原价的 $85\%$。
    $$ \text{原价} \times \text{折扣} = \text{现价} $$
    $$ \text{现价} \div \text{折扣} = \text{原价} $$
    $$ \text{现价} \div \text{原价} = \text{折扣} $$
\end{definition}

\begin{definition}[成数]
    表示一个数是另一个数的十分之几或百分之几十，通称“几成”。例如：二成就是（十分之二），改写成百分数是 $20\%$。
\end{definition}

\begin{definition}[税率]
    $$ \text{应纳税额} = \text{各种收入} \times \text{税率} $$
    $$ \text{各种收入} = \text{应纳税额} \div \text{税率} $$
\end{definition}

\begin{definition}[利率]
    存入银行的钱叫做本金。取款时银行多支付的钱叫做利息。利息和本金的比值叫做利率。
    $$ \text{利息} = \text{本金} \times \text{利率} \times \text{时间} $$
\end{definition}

\subsection*{题型}

\begin{example}
    王叔叔看中一套运动装，标价 200 元，经过还价，打八五折买到，王叔叔实际付了（ \quad ）元买了这套运动装。
\end{example}

\v{2em}

\begin{solution}
    $200 \times 85\% = 170$ (元)。
\end{solution}

\begin{example}
    一本书定价 75 元，售出后可获利 $50\%$，如果按定价的七折出售，可获利（ \quad ）元。
\end{example}

\v{2em}

\begin{solution}
    成本价: $75 \div (1+50\%) = 50$ (元)。
    七折售价: $75 \times 70\% = 52.5$ (元)。
    获利: $52.5 - 50 = 2.5$ (元)。
\end{solution}

\begin{example}
    王叔叔买了一辆价值 16000 元的摩托车。按规定，买摩托车要缴纳 $10\%$ 的车辆购置税。王叔叔买这辆摩托车一共要花多少钱？
\end{example}

\v{2em}

\begin{solution}
    购置税: $16000 \times 10\% = 1600$ (元)。
    总花费: $16000 + 1600 = 17600$ (元)。
\end{solution}

\begin{example}
    小强的妈妈在银行存了 5000 元，定期两年，年利率是 $4.50\%$，到期时，她应得利息（ \quad ）元。
\end{example}

\v{2em}

\begin{solution}
    利息: $5000 \times 4.50\% \times 2 = 450$ (元)。
\end{solution}


\section*{第三单元 圆柱和圆锥}

\subsection*{（一）圆柱}
\begin{definition}[圆柱的特征]
    \begin{enumerate}
        \item \textbf{底面的特征}：圆柱的底面是完全相同的两个圆。
        \item \textbf{侧面的特征}：圆柱的侧面是一个曲面。
        \item \textbf{高的特征}：圆柱有无数条高。
    \end{enumerate}
\end{definition}

\begin{definition}[圆柱的高]
    两个底面之间的距离叫做高。
\end{definition}

\begin{definition}[圆柱的侧面展开图]
    当沿高展开时展开图是长方形；当底面周长和高相等时，沿高展开图是正方形。
\end{definition}

\begin{theorem}[圆柱的侧面积]
    圆柱的侧面积 = 底面的周长 $\times$ 高，用字母表示为：
    $$ S_{侧} = C \times h \quad \text{或} \quad S_{侧} = 2\pi r \times h $$
\end{theorem}

\begin{theorem}[圆柱的表面积]
    圆柱的表面积 = 侧面积 + 2 个底面面积。
    $$ S_{表} = S_{侧} + 2 \times S_{底} \quad \text{或} \quad S_{表} = 2\pi r h + 2\pi r^2 $$
\end{theorem}

\begin{theorem}[圆柱的体积]
    圆柱所占空间的大小，叫做这个圆柱体的体积。
    $$ V = S_{底} \times h \quad \text{或} \quad V = \pi r^2 h $$
\end{theorem}

\subsection*{（二）圆锥}
\begin{definition}[圆锥]
    以直角三角形的一条直角边所在直线为旋转轴，其余两边旋转形成的旋转体叫做圆锥。该直角边叫圆锥的轴。
\end{definition}

\begin{definition}[圆锥的高]
    从圆锥的顶点到底面圆心的距离是圆锥的高。
\end{definition}

\begin{definition}[圆锥的特征]
    \begin{enumerate}
        \item \textbf{底面的特征}：圆锥的底面一个圆。
        \item \textbf{侧面的特征}：圆锥的侧面是一个曲面。
        \item \textbf{高的特征}：圆锥有一条高。
    \end{enumerate}
\end{definition}

\begin{theorem}[圆锥的侧面展开图]
    把圆锥的侧面展开得到一个扇形。
\end{theorem}

\begin{theorem}[圆锥的体积]
    一个圆锥所占空间的大小，叫做这个圆锥的体积。一个圆锥的体积等于与它等底等高的圆柱的体积的 $\frac{1}{3}$。
    $$ V = \frac{1}{3} S_{底} \times h \quad \text{或} \quad V = \frac{1}{3} \pi r^2 h $$
\end{theorem}

\begin{theorem}[圆柱与圆锥的关系]
    \begin{enumerate}
        \item 与圆柱等底等高的圆锥体积是圆柱体积的三分之一。
        \item 体积和高相等的圆锥与圆柱之间，圆锥的底面积是圆柱的三倍。
        \item 体积和底面积相等的圆锥与圆柱之间，圆锥的高是圆柱的三倍。
    \end{enumerate}
\end{theorem}

\subsection*{题型}
\begin{example}
    一个圆柱的底面半径是 5cm，高是 10cm，它的底面积是（ \quad ）$cm^2$，侧面积是（ \quad ）$cm^2$，体积是（ \quad ）$cm^3$。
\end{example}

\v{2em}

\begin{solution}
    底面积 $S_{底} = \pi r^2 = \pi \times 5^2 = 25\pi \ (cm^2)$。\\
    侧面积 $S_{侧} = 2\pi r h = 2\pi \times 5 \times 10 = 100\pi \ (cm^2)$。\\
    体积 $V = S_{底} \times h = 25\pi \times 10 = 250\pi \ (cm^3)$。
\end{solution}

\begin{example}
    用一张长 4.5 分米，宽 1.2 分米的长方形铁皮制成一个圆柱，这个圆柱的侧面积最多是（ \quad ）平方分米。（接口处不计）
\end{example}

\v{2em}

\begin{solution}
    侧面积就是这张长方形铁皮的面积：$4.5 \times 1.2 = 5.4$ (平方分米)。
\end{solution}

\begin{example}
    一个圆锥和一个圆柱等底等高，圆锥的体积是 $76cm^3$，圆柱的体积是（ \quad ）$cm^3$。
\end{example}

\v{2em}

\begin{solution}
    $V_{圆柱} = 3 \times V_{圆锥} = 3 \times 76 = 228 \ (cm^3)$。
\end{solution}

\begin{example}
    一个圆锥的底面直径和高都是 6cm, 它的体积是( \quad ) $cm^3$。
\end{example}

\v{2em}

\begin{solution}
    半径 $r = 6 \div 2 = 3$ (cm)。
    体积 $V = \frac{1}{3}\pi r^2 h = \frac{1}{3} \pi \times 3^2 \times 6 = 18\pi \ (cm^3)$。
\end{solution}

\begin{example}
    求下面图形的体积。（单位：厘米）
\end{example}

\v{2em}

\begin{warning}
    缺少图形，无法解答。
\end{warning}

\begin{example}
    如图，先将甲容器注满水，再将水倒入乙容器，这时乙容器中的水有多高？（单位：厘米）
\end{example}

\v{2em}

\begin{warning}
    缺少图形，无法解答。
\end{warning}
\begin{tips}[解题思路]
    水的体积不变。设甲容器的底面半径为 $r_甲$、高为 $h_甲$，乙容器的底面半径为 $r_乙$。
    水的体积 $V = V_{甲} = \pi r_甲^2 h_甲$。
    将水倒入乙容器后，水的体积在乙容器中可以表示为 $V = \pi r_乙^2 h_{乙水面}$。
    因此 $h_{乙水面} = \frac{\pi r_甲^2 h_甲}{\pi r_乙^2} = \frac{r_甲^2 h_甲}{r_乙^2}$。
\end{tips}

\section*{第四单元 比例}

\subsection*{（一）比例的意义和基本性质}

\begin{definition}[比例的意义]
    表示两个比相等的式子叫做比例。如：$2:1 = 6:3$。
    组成比例的四个数，叫做比例的项。两端的两项叫做\textbf{外项}，中间的两项叫做\textbf{内项}。
\end{definition}

\begin{theorem}[比例的基本性质]
    在比例里，两个外项的积等于两个内项的积。这叫做比例的基本性质。
    例如：由 $3:2 = 6:4$ 可知 $3 \times 4 = 2 \times 6$；或者由 $x \times 1.5 = y \times 1.2$ 可知 $x:y = 1.2:1.5$。
\end{theorem}

\begin{tips}[比和比例的区别]
    \begin{enumerate}
        \item 比表示两个量相除的关系，它有两项（即前、后项）；比例表示两个比相等的式子，它有四项（即两个内项和两个外项）。
        \item 比有基本性质，它是化简比的依据；比例有基本性质，它是解比例的依据。
    \end{enumerate}
\end{tips}


\subsection*{（二）正比例和反比例}

\begin{definition}[成正比例的量]
    两种相关联的量，一种量变化，另一种量也随着变化，如果这两种量中相对应的两个数的比值（也就是商）一定，这两种量就叫做成正比例的量，他们的关系叫做正比例关系。用字母表示 $y/x = k$（一定）。\\
    \textbf{例如}：
    \begin{enumerate}
        \item 速度一定，路程和时间成正比例；因为：路程 $\div$ 时间 = 速度（一定）。
        \item 圆的周长和直径成正比例，因为：圆的周长 $\div$ 直径 = 圆周率（一定）。
        \item 圆的面积和半径不成比例，因为：圆的面积 $\div$ 半径 = 圆周率和半径的积（不一定）。
        \item $y = 5x$，y 和 x 成正比例，因为：$y \div x = 5$（一定）。
        \item 每天看的页数一定，总页数和天数成正比例，因为：总页数 $\div$ 天数 = 每天看页数（一定）。
    \end{enumerate}
\end{definition}

\begin{definition}[成反比例的量]
    两种相关联的量，一种量变化，另一种量也随着变化，如果这两种量中相对应的两个数的积一定，这两种量就叫做成反比例的量，他们的关系叫做反比例关系。用字母表示 $x \times y = k$（一定）。\\
    \textbf{例如}：
    \begin{enumerate}
        \item 路程一定，速度和时间成反比例，因为：速度 $\times$ 时间 = 路程（一定）。
        \item 总价一定，单价和数量成反比例，因为：单价 $\times$ 数量 = 总价（一定）。
        \item 长方形面积一定，它的长和宽成反比例，因为：长 $\times$ 宽 = 长方形的面积（一定）。
        \item $40 \div x = y$，x 和 y 成反比例，因为：$x \times y = 40$（一定）。
        \item 煤的总量一定，每天的烧煤量和烧的天数成反比例，因为：每天烧煤量 $\times$ 天数 = 煤的总量（一定）。
    \end{enumerate}
\end{definition}

\subsection*{（三）比例的应用}

\begin{definition}[比例尺]
    一幅图的图上距离和实际距离的比，叫做这幅图的比例尺。
    $$ \text{图上距离} : \text{实际距离} = \text{比例尺} $$
\end{definition}

\begin{tips}[比例尺的分类]
    \begin{enumerate}
        \item 数值比例尺和线段比例尺
        \item 缩小比例尺和放大比例尺
    \end{enumerate}
\end{tips}

\begin{theorem}[比例尺相关计算]
    $$ \text{实际距离} \times \text{比例尺} = \text{图上距离} $$
    $$ \text{图上距离} \div \text{比例尺} = \text{实际距离} $$
\end{theorem}


\begin{definition}[图形的放大与缩小]
    形状相同，大小不同。
\end{definition}

\subsection*{题型}
\begin{example}
    在一个比例中，两个内项正好互为倒数，已知一个外项是 $\frac{2}{5}$，则另一个外项是（ \quad ）。
\end{example}

\v{2em}

\begin{solution}
    两个内项互为倒数，则内项积为 1。根据比例的基本性质，外项积也为 1。所以另一个外项是 $\frac{5}{2}$。
\end{solution}

\begin{example}
    北京到天津的实际距离是 120 千米，在比例尺是 $\frac{1}{5000000}$ 的地图上，两地的图上距离是（ \quad ）厘米。
\end{example}

\v{2em}

\begin{solution}
    $120$ km $= 120,000$ m $= 12,000,000$ cm。\\
    图上距离 = $12,000,000 \times \frac{1}{5000000} = 2.4$ (厘米)。
\end{solution}

\begin{example}
    如果 $2a=3b$，那么 $a:b=$（ \quad ）:（ \quad ）。
\end{example}

\v{2em}

\begin{solution}
    由 $2a=3b$ 可得 $\frac{a}{b} = \frac{3}{2}$，所以 $a:b = 3:2$。
\end{solution}

\begin{example}
    在一副平面图上，用图上距离 2cm 表示实际距离 200m, 这幅图的比例尺是（ \quad ）\\
    A、 1:100 \qquad B、 1:1000 \qquad C 1:10000
\end{example}

\v{2em}

\begin{solution}
    200m = 20000cm。比例尺为 $2:20000 = 1:10000$。选 C。
\end{solution}

\begin{example}
    按 1:5 将长方形缩小，就是将长方形的面积缩小到原来的（ \quad ）\\
    A、 $\frac{1}{5}$ \qquad B、 $\frac{1}{10}$ \qquad C、 $\frac{1}{25}$
\end{example}

\v{2em}

\begin{solution}
    长和宽都缩小到原来的 $\frac{1}{5}$，面积缩小到原来的 $(\frac{1}{5})^2 = \frac{1}{25}$。选 C。
\end{solution}

\begin{example}
    算一算，解比例
    $$ x:10 = \frac{1}{4}:\frac{1}{3} \qquad 0.4:x = 1.2:2 \qquad \frac{1.2}{2.4} = \frac{3}{x} $$
\end{example}

\v{2em}

\begin{solution}
    \begin{itemize}
        \item $x:10 = \frac{1}{4}:\frac{1}{3} \implies \frac{1}{3}x = 10 \times \frac{1}{4} \implies \frac{1}{3}x = \frac{5}{2} \implies x = \frac{15}{2} = 7.5$
        \item $0.4:x = 1.2:2 \implies 1.2x = 0.4 \times 2 \implies 1.2x = 0.8 \implies x = \frac{0.8}{1.2} = \frac{2}{3}$
        \item $\frac{1.2}{2.4} = \frac{3}{x} \implies 1.2x = 2.4 \times 3 \implies 1.2x = 7.2 \implies x = 6$
    \end{itemize}
\end{solution}


\section*{第五单元 数学广角-鸽巢问题}
\begin{theorem}[抽屉原理（一）]
    把多于 n 个的物体放到 n 个抽屉里，则至少有一个抽屉里的东西不少于两件。
\end{theorem}

\begin{theorem}[抽屉原理（二）]
    把多于 $mn$ (m 乘以 n) 个的物体放到 n 个抽屉里，则至少有一个抽屉里有不少于 $m+1$ 的物体。
\end{theorem}

\begin{tips}
    抽屉原理解题的关键是正确地判断什么是抽屉，什么是物体？
    $$ \text{物体数} \div \text{抽屉数} = \text{商} \dots\dots \text{余数} $$
    $$ \text{至少数} = \text{商} + 1 $$
\end{tips}

\subsection*{题型}
\begin{example}
    一个小组 13 个人，其中至少有（ \quad ）人是同一个月出生的。
\end{example}

\v{2em}

\begin{solution}
    物体数：13 (人)，抽屉数：12 (月)。\\
    $13 \div 12 = 1 \dots\dots 1$。
    至少数 = $1+1=2$ (人)。
\end{solution}

\begin{example}
    6 只鸽子飞回 5 个鸽舍，至少有（ \quad ）只鸽子要飞进同一个鸽舍里。
\end{example}

\v{2em}

\begin{solution}
    物体数：6 (鸽子)，抽屉数：5 (鸽舍)。\\
    $6 \div 5 = 1 \dots\dots 1$。
    至少数 = $1+1=2$ (只)。
\end{solution}

\begin{example}
    7只兔子要装进6个笼子，至少有（ \quad ）只兔子要装进同一个笼子里。\\
    A. 3 \qquad B. 2 \qquad C. 4 \qquad D. 5
\end{example}

\v{2em}

\begin{solution}
    物体数：7 (兔子)，抽屉数：6 (笼子)。\\
    $7 \div 6 = 1 \dots\dots 1$。
    至少数 = $1+1=2$ (只)。选 B。
\end{solution}

\begin{example}
    张阿姨给孩子买衣服，有红、黄、白三种颜色，但结果总是至少有两个孩子的颜色一样，她至少有（ \quad ）孩子。\\
    A. 2 \qquad B. 3 \qquad C. 4 \qquad D. 6
\end{example}

\v{2em}

\begin{solution}
    抽屉数：3 (颜色)。要保证至少有2个孩子颜色一样，根据抽屉原理（一），物体数（孩子）必须比抽屉数多。
    至少孩子数 = $3+1=4$ (个)。选 C。
\end{solution}

\begin{example}
    7个人住进5个房间，至少要有两个人住同一间房。为什么？（请你用图示的方法说明理由）
\end{example}

\v{2em}

\begin{solution}
    \textbf{理由}：这是典型的鸽巢问题。我们可以把5个房间看作5个“抽屉”，7个人看作7个“物体”。

    \textbf{图示说明}：
    假设每个房间先住1个人：
    \begin{itemize}
        \item 房间1：1人
        \item 房间2：1人
        \item 房间3：1人
        \item 房间4：1人
        \item 房间5：1人
    \end{itemize}
    这样一共安排了5个人，还剩下 $7-5=2$ 个人。这剩下的2个人无论住进哪个房间，都会导致那个房间里有不止1个人。因此，至少要有两个人住同一间房。
\end{solution}

\begin{example}
    把9本书放进2个抽屉里，总有一个抽屉至少放进5本书，为什么？
\end{example}

\v{2em}

\begin{solution}
    \textbf{理由}：把2个抽屉看作“抽屉”，9本书看作“物体”。
    根据抽屉原理，我们进行计算：
    $$ 9 \div 2 = 4 \dots\dots 1 $$
    商是4，余数是1。这说明，平均分配后，每个抽屉可以放4本书，还剩下1本。
    这剩下的1本书，无论放入哪个抽屉，都会使那个抽屉的书本数量变成 $4+1=5$ 本。
    所以，总有一个抽屉至少放进5本书。
\end{solution}
